\begin{frame}{자주 하는 실수 (10.3)}
  \small
  \begin{block}{실수 1: 사전학습 모델의 \textbf{정규화 통계}를
  무시}
    사전학습된 모델은 보통 특정 정규화(예: ImageNet의
    \textbf{채널별 평균·표준편차}로 입력 표준화)를 적용받은
    데이터로 학습됐다. 새 데이터에 이 정규화를 \textbf{빼먹거나
    다른 값}을 쓰면, 필터들이 ``본 적 없는 스케일의 입력''을
    받아 성능이 예상보다 훨씬 나빠진다. \textbf{모델이
    학습될 때 가정했던 데이터 분포를 추론 시점에도 정확히
    재현해야 한다} --- 사전학습 모델을 쓸 때는 반드시
    \textbf{학습 전처리 파이프라인(평균/표준편차, 이미지
    크기)을 문서에서 확인하고 그대로 맞추는 것}이 첫 단계.
  \end{block}
  \vskip0.5em
  \begin{block}{실수 2: \textbf{높은 학습률}로 미세조정하며
  사전학습 가중치를 망침}
    사전학습 가중치는 사전학습 손실 지형에서 \textbf{좋은
    지역}에 놓인 점. 큰 학습률(예: 0.01)로 시작하면 그 좋은
    지역 경계를 \textbf{한 발로 뛰어넘는} 셈 --- 범용 특징이
    파괴되어 ``전이한 모델''이 ``전이를 안 한 것''보다
    \textbf{나빠짐}. 미세조정은 스크래치 대비
    \textbf{10$\sim$100배 낮은} 학습률(예: 0.001 또는
    0.0001)로, 가능하면 ``순차 동결 해제'' 순서로 ---
    학습률은 ``그래디언트 한 걸음의 대담함''(Week 9)이고,
    좋은 시작점 위에서는 \textbf{작은 걸음이 정답}.
  \end{block}
\end{frame}
